\section{Применение МКВ в рамках будущего исследования}
\label{sec:future-research}

\subsection{Постановка задачи и сценарий использования}
\label{subsec:mvp-scenario}

Обозначим сценарий работы системы: три организации совместно обучают логистическую регрессию для прогнозирования ухода клиента на записях с одинаковым набором числовых признаков и бинарной меткой. Исходные данные не раскрываются другим организациям и координатору.

Продукт состоит из двух независимых блоков: модели МО и подключаемого МКВ-модуля защищенного исполнения. Модуль работает на трех серверах, по одному от каждой организации; координатор запускает задания без доступа к данным, а результат инференса получает только владелец запроса. Обучение проводится на синтетических данных, результаты сопоставляются с открытым исполнением той же модели.

Цель разработки MVP --- продемонстрировать возможность безопасной реализации совместного обучения и инференса через заменяемый МКВ-модуль: без раскрытия входов в рамках получестной модели при честном большинстве (см. раздел~\ref{sec:theoretical-foundations}). Демонстрация подтверждает работоспособность интеграции, но не доказывает безопасность протокола; максимальная точность, устойчивость к активному нарушителю, справедливость и промышленная нагрузка не проверяются. Логистическая регрессия и готовые библиотечные операции выбраны, чтобы не реализовывать криптографические процедуры самостоятельно: это снижает риск ошибок реализации.

\subsection{Архитектура продукта}
\label{subsec:product-architecture}

Продукт делится на блок МО (модель и порядок обучения) и МКВ-модуль защищенного исполнения; подготовка данных и координация --- обслуживающие функции. Это позволяет проверять модель в открытом режиме и подключать защищенное исполнение через согласованный интерфейс. Приоритет --- безопасная интеграция заданного набора операций, а не универсальная поддержка моделей.

Блок МО описывает логистическую регрессию, порядок признаков и правила обновления весов. Вход обучения --- матрица признаков и вектор меток, вход инференса --- запись без метки, выход --- оценка вероятности ухода клиента; класс определяет владелец запроса по заданному порогу, что исключает отдельное защищенное сравнение.

МКВ-модуль выполняет операции над защищенными значениями на трех серверах: описание модели публично, а входы, метки, веса и промежуточные результаты представлены долями. Предварительное вычисление открытого предсказания с последующим разделением результата не допускается. Основания выбора криптографических методов и средств защиты информации приведены в разделе~\ref{sec:cryptographic-foundations} и таблице~\ref{tab:mpc-primitives}.

Интерфейс блоков включает инициализацию, обучение, инференс и сохранение состояния. Публичная конфигурация задает версию модели, размерности, параметры квантования, число итераций и получателя результата; приватные данные передаются ссылками на защищенные объекты. Блок МО задает операции, но не управляет обменом и восстановлением значений: скрытие данных от модели означает отсутствие открытых значений у кода, описывающего вычисления. Границы доступа приведены в таблице~\ref{tab:mvp-architecture}.

{\small
\setlength{\tabcolsep}{3pt}
\renewcommand{\arraystretch}{1.2}
\setlength{\tablecontentwidth}{\dimexpr\textwidth-6\tabcolsep-4\arrayrulewidth\relax}
\begin{longtable}{|P{0.22\tablecontentwidth}|P{0.39\tablecontentwidth}|P{0.39\tablecontentwidth}|}
\caption{Компоненты продукта и границы раскрытия данных}\label{tab:mvp-architecture}\\
\hline
\multicolumn{1}{|H{0.22\tablecontentwidth}|}{\textbf{Компонент}} & \multicolumn{1}{H{0.39\tablecontentwidth}|}{\textbf{Функция}} & \multicolumn{1}{H{0.39\tablecontentwidth}|}{\textbf{Доступные данные}} \\
\hline
\endfirsthead
\multicolumn{3}{@{}l@{}}{\normalsize Продолжение таблицы~\thetable}\\
\hline
\multicolumn{1}{|H{0.22\tablecontentwidth}|}{\textbf{Компонент}} & \multicolumn{1}{H{0.39\tablecontentwidth}|}{\textbf{Функция}} & \multicolumn{1}{H{0.39\tablecontentwidth}|}{\textbf{Доступные данные}} \\
\hline
\endhead
Клиент организации & Проверка, квантование и разделение входа; передача долей. & Собственные записи и метки; публичная конфигурация. \\
\hline
Блок МО & Описание модели и правил обучения. & Публичное описание; ссылки на защищенные значения. \\
\hline
МКВ-модуль: три сервера & Защищенное обучение и инференс; хранение состояния. & Локальные доли; сообщения протокола. \\
\hline
Координатор & Запуск заданий и сбор статусов. & Публичные настройки и статусы; без записей, долей и весов. \\
\hline
Владелец запроса & Восстановление результата и применение порога. & Собственный запрос и его предсказание. \\
\hline
\end{longtable}
}

Взаимодействие компонентов показано на рисунках~\ref{fig:mvp-dataflow-train} и~\ref{fig:mvp-dataflow-infer}; при обучении оно происходит в следующем порядке:
\begin{enumerate}
\item[1.] Координатор формирует задание с публичной конфигурацией (версия модели, размерности, параметры квантования, число итераций) и передает его клиентам организаций и МКВ-модулю.
\item[2.] Блок МО передает МКВ-модулю публичное описание модели и последовательность операций обучения.
\item[3.] Клиент каждой организации проверяет формат своих записей и меток, выполняет квантование, разделяет значения на доли и передает доли трем серверам МКВ-модуля; координатору записи и доли не передаются.
\item[4.] Серверы МКВ-модуля выполняют обновления весов над долями, обмениваясь сообщениями протокола; каждый сервер сохраняет собственную долю весов и идентификатор конфигурации.
\item[5.] Серверы сообщают координатору только статусы выполнения; координатор веса не восстанавливает.
\end{enumerate}

% Требуется в преамбуле: \usepackage{tikz}
\usetikzlibrary{arrows.meta}
\tikzset{
  mvpbox/.style={draw, rounded corners=2pt, minimum height=9mm, align=center, font=\footnotesize, fill=white, inner sep=3pt},
  mvporg/.style={mvpbox, fill=gray!12},
  mvpsrv/.style={mvpbox, fill=gray!25, minimum width=22mm},
  mvpmod/.style={draw, thick, rounded corners=4pt},
  mvpdata/.style={-{Stealth[length=2.4mm]}, very thick},
  mvpctrl/.style={-{Stealth[length=2.4mm]}, thick, dashed},
  mvpprot/.style={{Stealth[length=2mm]}-{Stealth[length=2mm]}, thick, densely dotted},
  mvpnum/.style={circle, draw, solid, fill=white, inner sep=0.6pt, minimum size=4.2mm, font=\scriptsize\bfseries},
  mvplab/.style={font=\scriptsize, align=center, inner sep=1pt},
}

\begin{figure}[!ht]
\centering
\resizebox{\ifdim\width>\linewidth\linewidth\else\width\fi}{!}{%
\begin{tikzpicture}
% МКВ-модуль и серверы
\node[mvpmod, minimum width=90mm, minimum height=36mm] (mod) at (10.3,0.7) {};
\node[font=\footnotesize\bfseries, anchor=north west] at ([xshift=3pt,yshift=-3pt]mod.north west) {МКВ-модуль};
\node[mvpsrv] (s1) at (7.3,1.25) {Сервер 1};
\node[mvpsrv] (s2) at (10.3,1.25) {Сервер 2};
\node[mvpsrv] (s3) at (13.3,1.25) {Сервер 3};
% остальные компоненты
\node[mvporg, minimum width=26mm, minimum height=14mm] (cl) at (1.2,1.25) {Клиенты\\организаций 1--3};
\node[mvpbox, minimum width=28mm] (co) at (10.3,4.1) {Координатор};
\node[mvpbox, minimum width=22mm] (mo) at (13.8,4.1) {Блок МО};
% шаг 1: задание
\draw[mvpctrl] (co.west) -| (cl.north);
\node[mvpnum] at (8.0,4.1) {1};
\node[mvplab, above] at (4.6,4.1) {задание и конфигурация};
\draw[mvpctrl] (9.6,3.65) -- (9.6,2.5) node[midway, mvpnum] {1} node[midway, left=3.5mm, mvplab] {задание};
% шаг 2: описание модели
\draw[mvpctrl] (13.8,3.65) -- (13.8,2.5) node[midway, mvpnum] {2} node[midway, right=3.5mm, mvplab] {описание\\модели};
% шаг 3: доли входов
\draw[mvpdata] (cl.east) -- (5.8,1.25) node[midway, mvpnum] {3} node[midway, below=2.5mm, mvplab] {доли записей\\и меток};
% шаг 4: протокол между серверами
\draw[mvpprot] (s1.east) -- (s2.west);
\draw[mvpprot] (s2.east) -- (s3.west);
\draw[mvpprot] (s1.south) to[bend right=20] (s3.south);
\node[mvpnum] at (6.3,-0.55) {4};
\node[mvplab, anchor=west] at (6.6,-0.55) {обновление весов над долями; каждый сервер хранит\\свою долю весов и идентификатор конфигурации};
% шаг 5: статусы
\draw[mvpctrl] (11.0,2.5) -- (11.0,3.65) node[midway, mvpnum] {5} node[midway, right=3.5mm, mvplab] {статусы};
\end{tikzpicture}}

\medskip
\resizebox{\ifdim\width>\linewidth\linewidth\else\width\fi}{!}{%
\begin{tikzpicture}
\draw[mvpdata] (0,0) -- (1.0,0);
\node[mvplab, anchor=west] at (1.1,0) {доли (приватные данные)};
\draw[mvpctrl] (5.6,0) -- (6.6,0);
\node[mvplab, anchor=west] at (6.7,0) {публичные данные и управление};
\draw[mvpprot] (11.8,0) -- (12.8,0);
\node[mvplab, anchor=west] at (12.9,0) {сообщения протокола};
\end{tikzpicture}}
\caption{Потоки данных между компонентами MVP при обучении}
\label{fig:mvp-dataflow-train}
\end{figure}

При инференсе компоненты взаимодействуют следующим образом:
\begin{enumerate}
\item[1.] Владелец запроса через клиент своей организации квантует запись, разделяет ее на доли и передает серверам МКВ-модуля; координатор запускает задание и получает только статусы.
\item[2.] Серверы загружают сохраненные доли весов, а блок МО задает операции инференса.
\item[3.] Серверы вычисляют оценку вероятности над долями и передают долю результата только владельцу запроса.
\item[4.] Владелец запроса восстанавливает вероятность и локально применяет порог классификации.
\end{enumerate}

\begin{figure}[!ht]
\centering
\resizebox{\ifdim\width>\linewidth\linewidth\else\width\fi}{!}{%
\begin{tikzpicture}
\node[mvpmod, minimum width=90mm, minimum height=36mm] (mod) at (10.3,0.7) {};
\node[font=\footnotesize\bfseries, anchor=north west] at ([xshift=3pt,yshift=-3pt]mod.north west) {МКВ-модуль};
\node[mvpsrv] (s1) at (7.3,1.25) {Сервер 1};
\node[mvpsrv] (s2) at (10.3,1.25) {Сервер 2};
\node[mvpsrv] (s3) at (13.3,1.25) {Сервер 3};
\node[mvporg, minimum width=38mm, minimum height=14mm] (ow) at (1.7,1.25) {Владелец запроса\\(клиент организации)};
\node[mvpbox, minimum width=28mm] (co) at (10.3,4.1) {Координатор};
\node[mvpbox, minimum width=22mm] (mo) at (13.8,4.1) {Блок МО};
% шаг 1: запуск и доли запроса
\draw[mvpctrl] (9.6,3.65) -- (9.6,2.5) node[midway, mvpnum] {1} node[midway, left=3.5mm, mvplab] {запуск\\задания};
\draw[mvpdata] (3.6,1.65) -- (5.8,1.65) node[midway, mvpnum] {1} node[midway, above=2.5mm, mvplab] {доли запроса};
\draw[mvpctrl] (11.0,2.5) -- (11.0,3.65) node[midway, right=3.5mm, mvplab] {статусы};
% шаг 2: операции инференса
\draw[mvpctrl] (13.8,3.65) -- (13.8,2.5) node[midway, mvpnum] {2} node[midway, right=3.5mm, mvplab] {операции\\инференса};
% шаг 3: результат
\draw[mvpdata] (5.8,0.85) -- (3.6,0.85) node[midway, mvpnum] {3} node[midway, below=2.5mm, mvplab] {доли результата};
% протокол между серверами
\draw[mvpprot] (s1.east) -- (s2.west);
\draw[mvpprot] (s2.east) -- (s3.west);
\draw[mvpprot] (s1.south) to[bend right=20] (s3.south);
\node[mvpnum] at (6.3,-0.3) {2};
\node[mvplab, anchor=west] at (6.6,-0.3) {загрузка сохраненных долей весов};
\node[mvpnum] at (6.3,-0.8) {3};
\node[mvplab, anchor=west] at (6.6,-0.8) {вычисление оценки вероятности над долями};
% шаг 4: восстановление у владельца
\node[mvpnum] at (0.1,-0.45) {4};
\node[mvplab, anchor=west] at (0.4,-0.45) {восстановление\\вероятности, порог};
\end{tikzpicture}}

\medskip
\resizebox{\ifdim\width>\linewidth\linewidth\else\width\fi}{!}{%
\begin{tikzpicture}
\draw[mvpdata] (0,0) -- (1.0,0);
\node[mvplab, anchor=west] at (1.1,0) {доли (приватные данные)};
\draw[mvpctrl] (5.6,0) -- (6.6,0);
\node[mvplab, anchor=west] at (6.7,0) {публичные данные и управление};
\draw[mvpprot] (11.8,0) -- (12.8,0);
\node[mvplab, anchor=west] at (12.9,0) {сообщения протокола};
\end{tikzpicture}}
\caption{Потоки данных между компонентами MVP при инференсе}
\label{fig:mvp-dataflow-infer}
\end{figure}

Для проверки блок МО независимо исполняется в открытом режиме на синтетических данных, а результаты сопоставляются с защищенным исполнением. Журналы всех компонентов содержат только статусы, размерности и время выполнения; скрытие служебных сведений в MVP не предусматривается. Порядок признаков, масштабирование и границы значений согласуются до запуска.

Заменяемость модуля требует совместимости интерфейса, числового представления и правил раскрытия. MVP использует адаптер с фиксированными операциями логистической регрессии; замена реализации требует проверки численной совместимости и модели угроз, а оптимизация обмена и расширение вычислений рассматриваются позже с учетом раздела~\ref{sec:limitations}.

\subsection{Модель угроз и границы MVP}
\label{subsec:mvp-threat-model}

Модель угроз конкретизирует постановку из подразделов~\ref{subsec:formal-statement}--\ref{subsec:adversary-models} применительно к архитектуре подраздела~\ref{subsec:product-architecture}: три организации, три сервера (по одному от каждой), координатор без доступа к приватным входам и долям. Допущения приведены в таблице~\ref{tab:mvp-threat-model}.

{\small
\setlength{\tabcolsep}{3pt}
\renewcommand{\arraystretch}{1.2}
\setlength{\tablecontentwidth}{\dimexpr\textwidth-4\tabcolsep-3\arrayrulewidth\relax}
\begin{longtable}{|P{0.27\tablecontentwidth}|P{0.73\tablecontentwidth}|}
\caption{Допущения модели угроз MVP}\label{tab:mvp-threat-model}\\
\hline
\multicolumn{1}{|H{0.27\tablecontentwidth}|}{\textbf{Параметр}} & \multicolumn{1}{H{0.73\tablecontentwidth}|}{\textbf{Принятое допущение}} \\
\hline
\endfirsthead
\multicolumn{2}{@{}l@{}}{\normalsize Продолжение таблицы~\thetable}\\
\hline
\multicolumn{1}{|H{0.27\tablecontentwidth}|}{\textbf{Параметр}} & \multicolumn{1}{H{0.73\tablecontentwidth}|}{\textbf{Принятое допущение}} \\
\hline
\endhead
Нарушитель & Получестный: наблюдает состояние своего сервера и сообщения, но не изменяет предписанные операции. \\
\hline
Допустимый сговор & Не более одной организации (ее клиент и сервер); сговор с координатором допускается. \\
\hline
Честное большинство & Не менее двух серверов честны; административные области независимы. \\
\hline
Каналы и среда & Аутентифицированные конфиденциальные каналы; корректная реализация и случайность. Компрометация среды и побочные каналы исключены. \\
\hline
Допустимое раскрытие & Конфигурация, размеры выборок, время и статусы; предсказание --- владельцу запроса. Веса и промежуточные значения не восстанавливаются. \\
\hline
\end{longtable}
}

Защите подлежат чужие записи, метки и промежуточные значения. Допустимость сговора --- требование к протоколу, а не следствие наличия трех серверов; три процесса на одном компьютере проверяют интеграцию, но не независимость административных областей. Зависимости от вспомогательных компонентов подготовки вычислений учитываются при выборе платформы.

Вне MVP остаются защита от активного нарушителя, проверка истинности данных, подмена долей, намеренное прекращение участия, устойчивость к отказам, гарантированная выдача результата и справедливость: при отказе сервера задание завершается с ошибкой. Вывод сведений об обучающей выборке из разрешенных предсказаний, DP и скрытие служебных сведений также не предусматриваются (см. раздел~\ref{sec:limitations}); гарантии за пределами принятых допущений не заявляются.

\subsection{Выбор протокола и платформы}
\label{subsec:mvp-protocol-platform}

Согласно таблице~\ref{tab:mpc-primitives}, арифметическая часть модели строится на разделении секрета. Базовое сочетание --- TFEncrypted и получестный вариант ABY3 для трех серверов \cite{20}\cite{21}. Выбор обусловлен доступностью библиотечных операций и соответствием допущениям подраздела~\ref{subsec:mvp-threat-model}. Обзор применений платформ приведен в подразделе~\ref{subsec:ml-applications}; их сопоставление для MVP --- в таблице~\ref{tab:mvp-platforms}.

{\small
\setlength{\tabcolsep}{3pt}
\renewcommand{\arraystretch}{1.2}
\setlength{\tablecontentwidth}{\dimexpr\textwidth-6\tabcolsep-4\arrayrulewidth\relax}
\begin{longtable}{|P{0.19\tablecontentwidth}|P{0.39\tablecontentwidth}|P{0.42\tablecontentwidth}|}
\caption{Сопоставление платформ для реализации MVP}\label{tab:mvp-platforms}\\
\hline
\multicolumn{1}{|H{0.19\tablecontentwidth}|}{\textbf{Платформа}} & \multicolumn{1}{H{0.39\tablecontentwidth}|}{\textbf{Основание применения}} & \multicolumn{1}{H{0.42\tablecontentwidth}|}{\textbf{Решение для MVP}} \\
\hline
\endfirsthead
\multicolumn{3}{@{}l@{}}{\normalsize Продолжение таблицы~\thetable}\\
\hline
\multicolumn{1}{|H{0.19\tablecontentwidth}|}{\textbf{Платформа}} & \multicolumn{1}{H{0.39\tablecontentwidth}|}{\textbf{Основание применения}} & \multicolumn{1}{H{0.42\tablecontentwidth}|}{\textbf{Решение для MVP}} \\
\hline
\endhead
CrypTen \cite{22} & Интерфейс на основе PyTorch; арифметические и бинарные доли. & Не выбран: режим с доверенным генератором троек добавляет допущение; репозиторий архивирован. \\
\hline
TFEncrypted \cite{21} & Интерфейс TensorFlow; ABY3 на трех серверах; обучение и инференс. & Выбран: готовые операции соответствуют составу участников; потребуется фиксация совместимых версий. \\
\hline
SecureML \cite{23} & Исследовательская реализация МО на двух не вступающих в сговор серверах. & Не выбран: перенос в трехсерверную архитектуру потребует дополнительных решений. \\
\hline
\end{longtable}
}

Линейные операции выполняются на реплицированном аддитивном разделении ABY3 без самостоятельной реализации схемы Шамира \cite{20}\cite{21}; сложение, умножение матриц и усечение чисел с фиксированной запятой выполняет библиотека. Для сигмоиды используется библиотечная аппроксимация с тремя линейными участками из ABY3 \cite{21}: выбор участка выполняется защищенными сравнениями над бинарными долями, отдельная подсистема GC не интегрируется, а используемый библиотекой OT не реализуется заново. Нелинейная обработка остается внутри МКВ-модуля.

Обучение проходит фиксированное число шагов без защищенного критерия остановки. Открытая контрольная реализация применяет ту же аппроксимацию; влияние аппроксимации на качество оценивается дополнительно относительно исходной сигмоиды. Пороговая классификация остается локальной операцией владельца результата согласно подразделу~\ref{subsec:product-architecture}.

FHE, самостоятельная интеграция HE и GC и протоколы защиты от активного нарушителя для MVP не выбираются, так как выходят за его границы и принятую получестную модель угроз (см. таблицу~\ref{tab:mpc-primitives}); это не означает непригодности указанных подходов.

Случайность готовится средствами выбранного трехстороннего протокола, без доверенного генератора троек и без передачи этой роли координатору \cite{20}\cite{21}. Для воспроизводимости фиксируется версия TFEncrypted; поскольку TFEncrypted позиционируется как экспериментальное программное обеспечение \cite{21}, его выбор не заменяет проверки реализации. Адаптер из подраздела~\ref{subsec:product-architecture} позволяет заменить платформу без изменения описания модели.

\subsection{Функциональные требования к системе}
\label{subsec:mvp-requirements}

Требования определяют проверяемый объем прототипа в рамках подразделов~\ref{subsec:mvp-scenario}--\ref{subsec:mvp-protocol-platform}. Функциональные требования приведены в таблице~\ref{tab:mvp-functional-requirements}.

{\small
\setlength{\tabcolsep}{3pt}
\renewcommand{\arraystretch}{1.2}
\setlength{\tablecontentwidth}{\dimexpr\textwidth-4\tabcolsep-3\arrayrulewidth\relax}
\begin{longtable}{|P{0.12\tablecontentwidth}|P{0.88\tablecontentwidth}|}
\caption{Функциональные требования к MVP}\label{tab:mvp-functional-requirements}\\
\hline
\multicolumn{1}{|H{0.12\tablecontentwidth}|}{\textbf{Код}} & \multicolumn{1}{H{0.88\tablecontentwidth}|}{\textbf{Требование}} \\
\hline
\endfirsthead
\multicolumn{2}{@{}l@{}}{\normalsize Продолжение таблицы~\thetable}\\
\hline
\multicolumn{1}{|H{0.12\tablecontentwidth}|}{\textbf{Код}} & \multicolumn{1}{H{0.88\tablecontentwidth}|}{\textbf{Требование}} \\
\hline
\endhead
Ф1 & Система должна принимать локальные файлы CSV и отклонять входы с неверной размерностью, диапазоном или пропусками до передачи данных. \\
\hline
Ф2 & Система должна выполнять квантование и разделение входов у владельца и передавать доли трем серверам; записи и доли координатору не передаются. \\
\hline
Ф3 & Система должна обучать логистическую регрессию через МКВ-модуль с заданным числом обновлений и библиотечной аппроксимацией сигмоиды. \\
\hline
Ф4 & Система должна сохранять и загружать состояние модели как локальные доли на серверах. \\
\hline
Ф5 & Система должна выполнять защищенный инференс и выдавать вероятность только владельцу запроса. \\
\hline
Ф6 & Система должна обеспечивать открытое контрольное исполнение и отчет о качестве, времени и трафике. \\
\hline
\end{longtable}
}


\subsection{План реализации и критерии готовности MVP}
\label{subsec:mvp-implementation-plan}

Этапы приведены в таблице~\ref{tab:mvp-implementation-stages}: используются готовые операции с интерфейсом подраздела~\ref{subsec:product-architecture} без собственных криптографических процедур. Метки синтетических данных задаются порогом линейной комбинации признаков, что упрощает проверку обучения, но не воспроизводит статистику реальных клиентов.

{\small
\setlength{\tabcolsep}{3pt}
\renewcommand{\arraystretch}{1.2}
\setlength{\tablecontentwidth}{\dimexpr\textwidth-6\tabcolsep-4\arrayrulewidth\relax}
\begin{longtable}{|P{0.22\tablecontentwidth}|P{0.40\tablecontentwidth}|P{0.38\tablecontentwidth}|}
\caption{Этапы реализации MVP и ожидаемые результаты}\label{tab:mvp-implementation-stages}\\
\hline
\multicolumn{1}{|H{0.22\tablecontentwidth}|}{\textbf{Этап}} & \multicolumn{1}{H{0.40\tablecontentwidth}|}{\textbf{Содержание работы}} & \multicolumn{1}{H{0.38\tablecontentwidth}|}{\textbf{Ожидаемый результат}} \\
\hline
\endfirsthead
\multicolumn{3}{@{}l@{}}{\normalsize Продолжение таблицы~\thetable}\\
\hline
\multicolumn{1}{|H{0.22\tablecontentwidth}|}{\textbf{Этап}} & \multicolumn{1}{H{0.40\tablecontentwidth}|}{\textbf{Содержание работы}} & \multicolumn{1}{H{0.38\tablecontentwidth}|}{\textbf{Ожидаемый результат}} \\
\hline
\endhead
Данные & Генерация выборок; запуск открытой модели. & Файлы обучения и тестовая выборка; контрольные предсказания. \\
\hline
Квантование & Преобразование входов в формат с фиксированной запятой. & Согласованный числовой формат. \\
\hline
Разделение на доли & Передача долей серверам. & Доли на серверах; записи остаются у владельцев. \\
\hline
Вычисление & 150 обновлений весов; сохранение состояния; инференс после перезапуска. & Загружаемые доли весов; ответы на тестовые записи. \\
\hline
Результат & Выдача вероятностей владельцу; сопоставление с открытым исполнением. & Отчет о качестве, времени и обмене. \\
\hline
\end{longtable}
}

MVP считается готовым при выполнении условий таблицы~\ref{tab:mvp-readiness}. Проверки проводятся на трех раздельных узлах; локальный запуск трех процессов допустим только при отладке, согласно ограничениям подраздела~\ref{subsec:mvp-threat-model}.

{\small
\setlength{\tabcolsep}{3pt}
\renewcommand{\arraystretch}{1.2}
\setlength{\tablecontentwidth}{\dimexpr\textwidth-4\tabcolsep-3\arrayrulewidth\relax}
\begin{longtable}{|P{0.25\tablecontentwidth}|P{0.75\tablecontentwidth}|}
\caption{Критерии готовности MVP}\label{tab:mvp-readiness}\\
\hline
\multicolumn{1}{|H{0.25\tablecontentwidth}|}{\textbf{Проверка}} & \multicolumn{1}{H{0.75\tablecontentwidth}|}{\textbf{Условие приемки}} \\
\hline
\endfirsthead
\multicolumn{2}{@{}l@{}}{\normalsize Продолжение таблицы~\thetable}\\
\hline
\multicolumn{1}{|H{0.25\tablecontentwidth}|}{\textbf{Проверка}} & \multicolumn{1}{H{0.75\tablecontentwidth}|}{\textbf{Условие приемки}} \\
\hline
\endhead
Запуск и входы & Сценарий запускает обучение и инференс на 900 и 300 записях; некорректные входы отклоняются. \\
\hline
Обучение и состояние & Выполнены 150 обновлений; после перезапуска инференс использует сохраненные доли; точность выше постоянного прогноза преобладающего класса. \\
\hline
Численная совместимость & В каждом из пяти запусков снижение точности не более 2 процентных пунктов, отклонение вероятностей --- не более $0{,}005$. \\
\hline
Границы раскрытия & Проверка интерфейсов и трассировки сообщений подтверждает отсутствие восстановления входов, весов и градиентов; журналы не содержат запрещенных значений из подраздела~\ref{subsec:mvp-requirements}. \\
\hline
Отчет & Для пяти запусков сохранены конфигурации, версии и метрики; трафик шага сопоставлен с пределом 10~МиБ. \\
\hline
\end{longtable}
}

Проверка интерфейсов и трассировки сообщений не заменяет доказательство безопасности; готовность определяется в рамках подраздела~\ref{subsec:mvp-threat-model}.

\subsection{Риски, ограничения и направления развития после MVP}
\label{subsec:mvp-risks-roadmap}

Ограничения прототипа определяются подразделом~\ref{subsec:mvp-threat-model} и выводами раздела~\ref{sec:limitations}. Направления таблицы~\ref{tab:mvp-risks-roadmap} не входят в критерии готовности MVP, так как выходят за принятые границы MVP и модель угроз. Риски несовместимости зависимостей и числового представления контролируются фиксацией версий и проверками подразделов~\ref{subsec:mvp-requirements}--\ref{subsec:mvp-implementation-plan}.

{\small
\setlength{\tabcolsep}{3pt}
\renewcommand{\arraystretch}{1.2}
\setlength{\tablecontentwidth}{\dimexpr\textwidth-6\tabcolsep-4\arrayrulewidth\relax}
\begin{longtable}{|P{0.23\tablecontentwidth}|P{0.34\tablecontentwidth}|P{0.43\tablecontentwidth}|}
\caption{Ограничения MVP и направления последующего развития}\label{tab:mvp-risks-roadmap}\\
\hline
\multicolumn{1}{|H{0.23\tablecontentwidth}|}{\textbf{Область риска}} & \multicolumn{1}{H{0.34\tablecontentwidth}|}{\textbf{Граница MVP}} & \multicolumn{1}{H{0.43\tablecontentwidth}|}{\textbf{Развитие после MVP}} \\
\hline
\endfirsthead
\multicolumn{3}{@{}l@{}}{\normalsize Продолжение таблицы~\thetable}\\
\hline
\multicolumn{1}{|H{0.23\tablecontentwidth}|}{\textbf{Область риска}} & \multicolumn{1}{H{0.34\tablecontentwidth}|}{\textbf{Граница MVP}} & \multicolumn{1}{H{0.43\tablecontentwidth}|}{\textbf{Развитие после MVP}} \\
\hline
\endhead
Точность и издержки & Синтетические данные и одна модель; промышленная нагрузка не проверяется. & Расширить испытания; оптимизировать объем обмена согласно подразделу~\ref{subsec:communication}. \\
\hline
Модель противника & Активные нарушения и сговор двух организаций вне модели; справедливость не обеспечивается. & Выбрать иной профиль безопасности и заменить реализацию МКВ-модуля. \\
\hline
Раскрытие через ответы & DP и защита от вывода сведений по повторным запросам не реализуются. & Исследовать DP внутри МКВ-модуля до выдачи ответа согласно подразделу~\ref{subsec:dp}. \\
\hline
Среда исполнения & Компрометация среды и побочные каналы исключены; TEE не используется. & Рассмотреть отдельный профиль с TEE, пересмотрев доверие к аппаратной среде согласно подразделу~\ref{subsec:tee}. \\
\hline
Сценарий обучения & Совместное обучение на долях; режим FL не предусмотрен. & Исследовать локальное обучение и защищенное агрегирование обновлений; отдельно определить доступ к весам согласно подразделу~\ref{subsec:dp}. \\
\hline
Служебные сведения & Размерности и время не скрываются; ORAM не применяется. & Оценить допустимость утечек; при необходимости исследовать ORAM согласно подразделу~\ref{subsec:oram}. \\
\hline
\end{longtable}
}

Первым этапом после приемки станут испытания на дополнительных данных и анализ измеренных издержек, затем --- изменения протокола и механизма выдачи ответа после уточнения модели угроз. TEE, FL и ORAM вводятся отдельными экспериментами, чтобы определить вклад каждого изменения. Граница блоков из подраздела~\ref{subsec:product-architecture} позволяет сохранять описание модели при замене защищенных операций, однако каждая замена требует повторных проверок численной совместимости и безопасности, а практическое внедрение --- организационных процедур подраздела~\ref{subsec:standardization}.

\subsection{Вывод}
\label{subsec:mvp-conclusion}

Предложен план разработки продукта из двух раздельных блоков: модели МО и заменяемого МКВ-модуля. Цель реализации --- показать возможность безопасного совместного обучения и инференса логистической регрессии на данных трех организаций без раскрытия входов другим участникам и координатору при допущениях подраздела~\ref{subsec:mvp-threat-model}. Выполнение критериев подраздела~\ref{subsec:mvp-implementation-plan} подтвердит работоспособность интеграции и численные допуски, но не заменит доказательство безопасности и не означает готовности к промышленному внедрению. План связывает анализ основ, классификации и применений МКВ с проверяемым экспериментом; разделение блоков позволяет развивать продукт по направлениям подраздела~\ref{subsec:mvp-risks-roadmap} без включения криптографической логики в описание модели.